\begin{textsample}{NEG Dim 1 – human – Score -6.00 – rj/rj\_ef\_linguagens\_ed\_fis\_1\_p9.txt}  \label{ex:f1_neg_010}
Dessa \textbf{forma}, a Educação Física Escolar \textbf{tem} como objetivo, mais do que apresentar -se como fonte de conhecimentos, despertar nos estudantes o interesse em envolver -se com as atividades e exercícios \textbf{corporais}, \textbf{possibilitando} convivências harmoniosas e construtivas com outros indivíduos, sendo eles capazes de reconhecer e \textbf{respeitar} as características físicas e desempenho de si próprio e de outros indivíduos, não segregando e nem depreciando outras pessoas por qualidades e peculiaridades como aspectos físicos, sexuais e ou sociais.
Com a Educação Física escolar é possível evidenciar a liberdade cognitiva e emocional dos estudantes para a aprendizagem e o saber como relacionar -se em \textbf{grupo} por meio da evolução positiva de comportamentos, valores, normas e atitudes.

% matched lemmas: corporal, forma, grupo, possibilitar, respeitar, ter
\end{textsample}
